\chapter{Introduction}

% \section{Background}
% \begin{itemize}
%     \item Evolution of millimeter-wave communication systems
%     \begin{itemize}
%         \item FR1 (cellular bands becoming saturated, motivating the shift to mmWave)
%     \end{itemize}
%     \item Challenges of line-of-sight-dominant MIMO channels
%     \begin{itemize}
%         \item Discuss common mmWave channel challenges, supported by literature
%     \end{itemize}
%     \item Importance of channel decorrelation and spatial multiplexing
%     \item Motivation for antenna lens-assisted MIMO systems
%     \begin{itemize}
%         \item Improve channel decorrelation and spatial multiplexing
%     \end{itemize}
% \end{itemize}


\section{Background}

% \paragraph{1 - Development of 5G and the Utilization of FR2 Spectrum}
% Paragraf pertama membuka pembahasan dari kebutuhan komunikasi nirkabel berkecepatan tinggi dan penggunaan spektrum frekuensi tinggi pada 5G New Radio.

% Poin utama:

% Pertumbuhan kebutuhan data dan perangkat terhubung membutuhkan bandwidth komunikasi yang lebih luas karena scarcity and fragmentation in conventional cellular bands dan pindah ke mmwave\cite{Rappaport2013}
% Salah satu pengembangan penting dalam 5G New Radio adalah pemanfaatan Frequency Range 2 (FR2) yang memiliki bandwidth yang luas \cite{TSGR2021}.
% Frekuensi penelitian sebesar 38 GHz berada dalam FR2-1.
% Secara lebih spesifik, 38 GHz termasuk dalam NR operating band n260, yang mencakup 37–40 GHz yang menyediakan hampir 3 GHz bandwidth untuk komunikasi nirkabel berkecepatan tinggi.
% % Sitasi : \cite{Rappaport2013,TSGR2021}
The rapid growth in data traffic and the proliferation of connected devices continue to increase the bandwidth required by next-generation wireless networks. Because conventional cellular spectrum is scarce and fragmented, millimeter-wave (mmWave) frequencies have become an important means of providing wider transmission bandwidths \cite{Rappaport2013}. Accordingly, 5G New Radio (NR) includes Frequency Range~2 (FR2) for high-capacity wireless communication \cite{TSGR2021}. This thesis focuses on 38~GHz, which lies within FR2-1 and NR operating band n260; the nominal frequency range of n260 is 37--40~GHz.

% \paragraph{Propagational Challenges and Hardware Complexity in 5G FR2}
% The utilization of the high-frequency FR2 spectrum is inherently linked to significant propagation challenges arising from its short wavelengths. Signals at these frequencies suffer from severe free-space path loss, substantial penetration losses, and extreme sensitivity to physical blockages, which collectively restrict overall coverage \cite{Zeng2016}. Consequently, maintaining reliable wireless links in FR2 environments requires sophisticated techniques such as directional transmissions and high antenna gains. While the small wavelength size facilitates the integration of numerous antenna elements into a compact aperture, it also necessitates advanced spatial processing capabilities. Conventional approaches to achieve this include Massive MIMO systems designed to provide sufficient beamforming gain; however, fully digital MIMO architectures—which require an independent RF chain for every antenna element—pose substantial practical hurdles at FR2 frequencies. The high-frequency environment demands complex and costly components such as power amplifiers (PAs), low-noise amplifiers (LNAs), mixers, local oscillators, and high-speed ADCs or DACs \cite{Heath2016}. Scaling the number of RF chains in these systems significantly escalates hardware complexity, power consumption, manufacturing costs, and the stringency of calibration requirements.

% Poin utama Paragraf 2:
% Sinyal FR2 memiliki panjang gelombang yang sangat pendek.
% Frekuensi tinggi mengalami:
% free-space path loss yang tinggi;
% penetration loss yang besar;
% sensitivitas terhadap blockage;
% cakupan propagasi yang lebih terbatas.
% Karena itu, sistem FR2 memerlukan directional transmission dan antenna gain yang tinggi.
% Ukuran panjang gelombang yang kecil memungkinkan banyak elemen antena ditempatkan pada aperture yang relatif ringkas. \cite{Zeng2016}

% Poin utama Paragraf 3:
% Massive MIMO digunakan untuk menghasilkan beamforming gain.
% Fully digital MIMO membutuhkan satu RF chain untuk setiap elemen antena.
% Pada frekuensi FR2, RF chain memerlukan komponen seperti:
% power amplifier;
% low-noise amplifier;
% mixer;
% local oscillator;
% ADC atau DAC.
% Penambahan jumlah RF chain meningkatkan:
% hardware complexity;
% konsumsi daya;
% biaya implementasi;
% kebutuhan kalibrasi.
% sitasi : \cite{Heath2016}

% \paragraph{Lens Antenna Array for Beamspace Communication in FR2 MIMO}
An antenna-lens array provides an alternative front-end architecture by combining an electromagnetic lens with antenna elements positioned in the focal region or on a focal surface. The spatially varying phase delay across the lens aperture can passively transform signals from the antenna domain into the beamspace domain \cite{Shen2019}. Because waves arriving from different directions are focused onto different subsets of antenna elements, only the branches that capture dominant energy need to be connected to RF chains. This property is attractive for FR2 MIMO systems, whose channels are often dominated by a limited number of propagation paths and are consequently sparse in beamspace \cite{Zeng2016}. Under suitable angular-separation conditions, beam selection can therefore reduce the required number of RF chains, while the lens may improve angular discrimination and channel decorrelation and thereby support effective rank, spatial multiplexing, and channel capacity \cite{Zeng2016}.

% Poin utama Paragraf 4:
% Lens antenna array terdiri dari:
% electromagnetic lens;
% antenna elements pada focal region atau focal surface.
% Lens memberikan phase delay berbeda di sepanjang aperture.
% Gelombang dari arah berbeda difokuskan pada elemen antena berbeda.
% Lens secara pasif melakukan spatial transformation dari antenna space menuju beamspace.
% Hanya antena yang menerima energi dominan yang perlu dihubungkan ke RF chain.

% Poin utama Paragraf 5:
% Kanal FR2 umumnya memiliki jumlah dominant propagation paths yang terbatas.
% Beamspace channel menjadi relatif sparse.
% Sebagian besar energi kanal terkonsentrasi pada beberapa beam atau antena.
% Beam selection dapat mengurangi jumlah RF chain.
% Lens juga berpotensi meningkatkan:
% angular separation;
% spatial channel decorrelation;
% effective channel rank;
% spatial multiplexing capability;
% channel capacity.\cite{Zeng2016}

% \paragraph{Metasurface-Based Planar Lens Solutions for High-Frequency Communication}
Although conventional dielectric lenses can provide strong focusing and high gain, their thickness, weight, curved surfaces, and impedance-mismatch reflections can hinder integration with planar wireless front ends \cite{Al-Joumayly2011}. Planar metasurface lenses address these limitations by using arrays of subwavelength unit cells as local spatial phase shifters. Varying the unit-cell geometry tailors the transmission magnitude, phase, and frequency response across the aperture. The resulting phase profile can transform a spherical wavefront into an approximately planar wavefront, or conversely focus an incident planar wave, while retaining a compact, low-profile form factor \cite{Al-Joumayly2011}.

% Poin utama Paragraf 6:
% Dielectric lens dapat memberikan focusing dan gain tinggi.
% Namun, dielectric lens dapat memiliki:
% profil tebal;
% bobot besar;
% curved surface;
% reflection loss akibat impedance mismatch;
% integrasi yang sulit dengan planar antenna array;
% proses fabrikasi yang relatif kompleks.
% Masalah ini mendorong pengembangan planar lens.

% Poin utama Paragraf 7:
% Metasurface terdiri dari kumpulan sub-wavelength unit cells.
% Setiap unit cell bertindak sebagai spatial phase shifter.
% Dimensi geometris unit cell menentukan:
% transmission magnitude;
% transmission phase;
% frequency response.
% Distribusi unit cell di seluruh aperture membentuk phase profile yang diperlukan.
% Metasurface dapat mengubah:
% spherical wavefront menjadi planar wavefront;
% planar wavefront menjadi focused wavefront.

% \paragraph{8 — Relevance of a 38 GHz Metasurface Lens to 5G FR2}

% Paragraf ini membuat hubungan langsung antara desainmu dan 5G FR2.

% Poin utama:

% Lens dirancang pada 38 GHz.
% Frekuensi tersebut berada dalam NR band n260.
% Teknologi ini relevan untuk:
% FR2 access points;
% base stations;
% fixed wireless access;
% indoor high-capacity communication;
% directional MIMO receivers.
% Metasurface dapat melakukan passive beam focusing tanpa memerlukan phase shifter aktif pada setiap elemen.

\section{Literature Review}
\label{sec:literature_review}

The development of millimeter-wave (mmWave) wireless systems has
motivated extensive research on high-gain antennas, large-scale
multiple-input multiple-output (MIMO) architectures, and spatial
beamforming techniques. Although conventional antenna arrays can
provide substantial beamforming and spatial multiplexing gains, their
implementation becomes increasingly complex as the number of antenna
elements and radio-frequency (RF) chains increases. Electromagnetic
lens antennas provide an alternative approach by manipulating the
incident wavefront directly in the electromagnetic domain. When
combined with antenna elements positioned in the focal region, the
lens can also transform angle-dependent propagation components into
spatially separated focal responses.

This section reviews conventional mmWave MIMO systems, antenna-lens
technologies for beam focusing, beamspace MIMO and lens-assisted
communication systems, and existing physical implementations of antenna
lenses. The remaining research gaps are subsequently identified, with
particular emphasis on the feasibility and performance of applying
metasurface unit-cell designs to electrically larger simulated
apertures and the relationship
between physical lens focusing and MIMO channel characteristics.


\subsection{Conventional Millimeter-Wave MIMO Systems}
\label{subsec:conventional_mmwave_mimo}

Millimeter-wave communication provides access to substantially wider
frequency bandwidths than conventional sub-6-GHz systems. These wide
bandwidths are attractive for supporting the high data rates required
by future wireless applications. Nevertheless, mmWave propagation is
affected by substantial path loss, blockage, limited diffraction, and
relatively sparse scattering. Therefore, reliable mmWave links
generally require highly directional transmission and reception
\cite{Guo2026,Payami2020}.

The short wavelength at mmWave frequencies enables a large number of
antenna elements to be accommodated within a relatively compact
physical area. Large antenna arrays can therefore be used to generate
narrow beams and compensate for propagation loss through array
gain. In addition, when multiple independent propagation modes are
available, MIMO processing can support the simultaneous transmission
of multiple data streams. However, the practical advantages of
large-scale mmWave MIMO depend strongly on the transceiver
architecture.

In a fully digital MIMO architecture, each antenna element is connected
to an independent RF chain. This arrangement provides extensive
flexibility for precoding, combining, interference suppression, and
spatial multiplexing. However, each RF chain requires components such
as mixers, amplifiers, filters, and data converters. As the number of
antenna elements increases, the resulting hardware cost, power
consumption, and implementation complexity can become prohibitive
\cite{Payami2020,Zeng2018}.

Analog beamforming reduces the number of required RF chains by
connecting one RF chain to multiple antenna elements through a
phase-shifting network. This architecture can provide high array gain
with lower hardware complexity than a fully digital system. However,
a single-RF-chain analog beamformer generally supports only one data
stream and therefore cannot fully exploit the spatial-multiplexing
capability of a MIMO channel.

Hybrid analog--digital beamforming has consequently been introduced as
a compromise between fully digital and analog architectures. A
limited number of RF chains performs digital signal processing, while
an analog network connects these RF chains to a larger number of
antenna elements. Hybrid beamforming can provide a practical trade-off
among antenna gain, spatial multiplexing, hardware cost, and energy
consumption \cite{Payami2020}.

Nevertheless, fully connected hybrid architectures may require a
large number of phase shifters, power dividers, combiners, and RF
interconnections. Practical phase shifters also have finite phase
resolution and introduce insertion loss. In addition, the design of
wideband hybrid precoding and combining becomes more complicated
because the same analog network must operate over the complete signal
bandwidth.

The practical complexity of hybrid beamforming is illustrated by
Payami \textit{et al.}, who developed a fully connected hybrid
beamforming platform consisting of a $16\times8$ antenna array with
128 antenna elements and two RF inputs at 26~GHz
\cite{Payami2020}. Although the system provides independent
control of the phase and amplitude associated with each RF path, its
implementation requires extensive RF interconnections, calibration
procedures, control interfaces, phase shifters, and attenuators.

Furthermore, the sparse propagation characteristics of many mmWave
environments can limit the number of sufficiently independent spatial
modes. In line-of-sight-dominant scenarios, closely spaced antennas
with similar orientations may observe highly correlated channel
responses. Therefore, increasing the number of conventional antenna
elements does not necessarily produce a proportional increase in the
number of usable spatial streams.

These limitations have motivated the investigation of alternative
antenna architectures capable of providing high gain and spatial
discrimination without requiring a dedicated active phase-shifting
network for every antenna element. One promising solution is the use
of electromagnetic antenna lenses.


\subsection{Antenna Lens Technologies for Beam Focusing}
\label{subsec:antenna_lens_technologies}

An electromagnetic lens controls the spatial phase distribution of an
incident wave across its aperture. During transmission, a feed antenna
located near the focal point produces a spherical wave whose
propagation distance varies across the lens surface. The lens
introduces different transmission phases at different aperture
positions to compensate for these path-length differences. As a
result, the transmitted spherical wavefront can be transformed into
an approximately planar or shaped wavefront.

Through electromagnetic reciprocity, an incident plane wave can
conversely be concentrated within a focal region behind the lens.
This reciprocal property enables an antenna lens to operate both as a
high-gain transmitting aperture and as a spatial focusing device at
the receiver.

Conventional electromagnetic lenses are commonly implemented using
bulk dielectric materials. Their curved surfaces provide different
amounts of phase accumulation across the aperture. Dielectric lenses
can provide high antenna gain and relatively wide bandwidth,
particularly at mmWave frequencies. However, they may be bulky,
heavy, difficult to integrate with planar circuits, and affected by
reflection and dielectric losses.

Planar lens technologies reduce the physical profile by implementing
the required phase distribution using spatially distributed
phase-shifting elements. Early planar lenses used receiving and
transmitting antenna elements connected through transmission lines
with different lengths. Although these structures can provide the
required phase delay, their interconnection networks may become
complex for large apertures.

Frequency-selective surfaces, transmitarrays, metamaterials, and
metasurfaces provide alternative planar implementations. A
transmitarray generally consists of a spatial feed and a transmission
surface containing numerous unit cells. Each unit cell provides a
specific transmission magnitude and transmission phase. By arranging
unit cells with different responses over the aperture, the surface can
collimate, redirect, or shape the transmitted wave
\cite{Song2026}.

Al-Joumayly and Behdad developed a planar microwave lens based on
subwavelength spatial phase shifters derived from miniaturized-element
frequency-selective surfaces (MEFSSs) \cite{Al-Joumayly2011}. In their
approach, each subwavelength unit cell operates as a local spatial
phase shifter. Unit cells with different transmission phases are
distributed over the aperture to approximate the desired lens phase
profile. The fabricated X-band prototypes demonstrated approximately
20\% bandwidth and relatively stable responses under oblique
incidence.

Metasurfaces extend this concept by using electrically small
scatterers to manipulate the phase, amplitude, and polarization of
electromagnetic waves. Compared with conventional bulk lenses,
metasurface lenses can provide lower profiles, shorter focal
distances, reduced weight, and improved compatibility with printed
circuit manufacturing.

The performance of a metasurface lens depends strongly on its
unit-cell library. Important unit-cell characteristics include the
available transmission-phase range, transmission magnitude,
operating bandwidth, polarization response, angular stability,
physical periodicity, and fabrication sensitivity.

A unit-cell library should ideally provide a sufficiently wide phase
range while maintaining high and relatively uniform transmission
magnitude. Limited phase coverage can produce phase errors across the
aperture. Significant differences in transmission magnitude between
unit-cell geometries can create an unintended amplitude distribution,
reduce aperture efficiency, and increase radiation sidelobes.

The physical size of the unit cell is also important. Smaller unit
cells provide finer spatial sampling of the required phase
distribution and allow smoother transitions between neighboring
elements. However, very small structures may become more sensitive to
fabrication tolerances, conductor losses, substrate variations, and
mutual coupling.

The angular response of the unit cell is particularly important for a
spatially fed lens. Unit cells are frequently characterized using
periodic boundary conditions and normal plane-wave incidence. In an
actual lens, however, the feed illuminates different parts of the
aperture at different local incident angles. Unit cells near the
aperture edges may therefore exhibit transmission responses that
differ from their normal-incidence characteristics.

Consequently, designing a practical metasurface lens requires a
compromise among transmission efficiency, phase coverage, spatial
resolution, angular stability, bandwidth, aperture size, and
fabrication complexity.

The reciprocal focusing property also provides an important
connection between antenna-lens technology and MIMO communications.
Waves arriving from different angular directions can be concentrated
at different locations within the focal region. If several antenna
elements are positioned within this region, the lens can generate
angle-dependent spatial responses without requiring a conventional
active phase-shifter network.


\subsection{Beamspace MIMO and Lens-Assisted Systems}
\label{subsec:beamspace_mimo}

In conventional MIMO systems, signals and channels are generally
represented in the antenna-element domain. Beamspace MIMO instead
represents the transmitted and received signals using a set of
spatial beams. This transformation can concentrate the energy of a
sparse mmWave channel into a limited number of dominant beam
components.

An electromagnetic lens can approximately perform this spatial
transformation in the RF domain. A lens antenna array generally
consists of an electromagnetic lens and multiple antenna elements
placed along a focal arc or focal surface. An incident wave with a
particular angle of arrival is focused onto a localized subset of
antenna elements. Conversely, exciting different focal-region
elements produces beams with different departure angles.

Zeng and Zhang derived an analytical response for a lens antenna array
and showed that the energy distribution across the focal-region
antennas follows a sinc-like response
\cite{Zeng2016}. The position of the maximum response is determined
by the angle of arrival or departure. Consequently, propagation paths
with sufficiently separated angular directions can produce distinct
energy distributions across the focal-plane antenna elements.

This property is attractive for mmWave communications because mmWave
channels commonly contain only a limited number of dominant
propagation paths. In the beamspace representation, most of the
channel energy may therefore be concentrated into a small number of
beams. Only the antenna elements associated with these dominant beams
need to be connected to RF chains, potentially reducing the hardware
requirements and signal-processing complexity.

By exploiting the angular separation of different propagation paths,
Zeng and Zhang proposed path-division multiplexing and orthogonal
path-division multiplexing \cite{Zeng2016}. Under favorable
angular conditions, different propagation paths can be treated as
approximately parallel channels. Multiple data streams can then be
transmitted through different paths using relatively simple per-path
processing.

For propagation paths with insufficient angular separation,
small-scale group-based MIMO processing can be used to reduce
inter-path interference. This allows the lens-assisted architecture to
operate under more general propagation conditions than the ideal
orthogonal-path case.

The lens-assisted concept was extended to multiuser and
full-dimensional systems by Zeng, Yang, and Zhang
\cite{Zeng2018}. In a full-dimensional lens array, antenna
elements are distributed over a focal surface rather than only along
a one-dimensional focal arc. The arrangement therefore provides
angular resolution in both azimuth and elevation.

For wideband multiuser systems, path-delay compensation was introduced
to simplify the frequency-selective channels into smaller
approximately frequency-flat MIMO channels. The proposed architecture
achieved spectral efficiency comparable to conventional benchmark
systems while requiring fewer RF chains and lower signal-processing
complexity \cite{Zeng2018}.

Nevertheless, wideband lens-assisted MIMO systems are affected by beam
squint. The focusing location associated with a physical path can
change with frequency. Therefore, an antenna or beam selected at the
center frequency may not capture the complete channel energy across
the entire signal bandwidth.

Shen \textit{et al.} investigated this limitation and proposed a
phase-shifter-aided beam-selection network
\cite{Shen2019}. In the proposed architecture, a single RF
chain can support multiple focused-energy beams. The system combines
beam selection and beamspace precoding to capture the
frequency-dependent energy distribution without requiring an RF
chain for every beam.

These communication-oriented studies demonstrate the potential of
lens antenna arrays for beam selection, RF-chain reduction, and
path-based spatial multiplexing. However, most of
them model the lens using an ideal spatial transformation or an
analytical sinc-like response.

Such models provide valuable theoretical insight, but they do not
always include the full physical behavior of a practical
metasurface lens. These characteristics include finite aperture size,
insertion loss, focal-spot width, sidelobes, phase errors, unit-cell
coupling, manufacturing tolerances, and the actual radiation patterns
observed at different focal-plane antenna positions.


\subsection{Existing Research on Antenna Lenses}
\label{subsec:existing_antenna_lens_research}

Existing antenna-lens research can be broadly divided into the
physical development of electromagnetic lenses and
communication-oriented studies of lens-assisted systems. Physical
studies generally focus on
transmission magnitude, transmission phase, focusing intensity, focal
distance, antenna gain, bandwidth, polarization, and beam steering.
Communication-oriented studies instead emphasize beam selection,
spectral efficiency, path separation, RF-chain reduction, and
system-level throughput.

At microwave frequencies, Al-Joumayly and Behdad demonstrated
wideband planar lenses based on subwavelength spatial phase shifters
\cite{Al-Joumayly2011}. Their results showed that compact
periodic structures could provide different local transmission phases
while maintaining relatively stable responses over a useful bandwidth
and angular range. However, the implementation used multiple closely
spaced periodic layers to obtain the required transmission response.

At 28~GHz, several metasurface-lens functions have been demonstrated.
Yoon and Oh proposed a polarization-dependent thin lens based on
multiple rectangular patches and slotted grids
\cite{Yoon2019}. The lens produced different
electromagnetic responses for two orthogonal incident polarizations.
For one polarization, the structure operated as a focusing lens,
whereas for the other polarization, it behaved primarily as a
frequency-selective surface.

Tiwari \textit{et al.} developed a $25\times25$ digital metasurface
lens at 28~GHz with four discrete phase states
\cite{Tiwari2023}. Three switchable feed antennas were placed
below the lens to produce beams at broadside and approximately
$\pm30^{\circ}$. The measured broadside gain was 19~dBi, with a
reported gain enhancement of 14.9~dB and a scan loss of approximately
1~dB.

Although the approach provides active beam steering with relatively
few switching elements, the available beam directions are determined
by the finite number of feed positions. Increasing the number of beam
directions would require additional feed antennas, switching
components, or a more complex reconfigurable surface.

Lee \textit{et al.} proposed a large-aperture metamaterial lens for
multilayer MIMO transmission at 28~GHz
\cite{Lee2022Large}. The lens contained a $28\times28$
unit-cell array within an aperture of approximately
$50.4\times50.4$~mm$^2$. A channel model was developed by
representing propagation through the feed, metasurface, and
free-space regions as a sequence of channel components.

The measured and simulated results indicated lens gain of up to
approximately 14~dB, while system-level simulations demonstrated an
increase in user throughput. This study provides an important
connection between a physical metamaterial lens and communication
performance. However, its primary system-level objective was
throughput enhancement rather than a detailed investigation of
channel correlation and spatial-mode structure.

At frequencies around 38~GHz, Lee \textit{et al.} demonstrated a
single-layer phase-gradient metasurface lens operating from 35 to
40~GHz \cite{Lee2021}. The metasurface contained a
$13\times13$ array of unit cells with a period of 3~mm, resulting in
an aperture size of $39\times39$~mm$^2$. The structure consisted of
two patterned metallic layers separated by a single dielectric
substrate.

By varying the geometrical parameters of the unit cells, the design
provided approximately full transmission-phase coverage while
maintaining satisfactory transmission magnitude. The metasurface was
designed for a nominal focal distance of 13~mm at 38~GHz and
demonstrated focusing over the frequency range from 35 to 40~GHz.

The focusing response was also characterized under different incident
angles. The lens maintained a spatially concentrated focal response
for incident angles up to approximately $30^{\circ}$. This result
demonstrates that a compact single-layer metasurface can provide
broadband and relatively angle-stable focusing around 38~GHz.

Nevertheless, the demonstrated lens was limited to a
$13\times13$ aperture. The study primarily evaluated the focal-field
distribution, focal position, transmission response, polarization,
and incident-angle stability. It did not investigate whether the same
unit-cell library could be extended to a substantially larger
aperture while maintaining the predicted phase and transmission
responses.

At higher mmWave frequencies, Bagheri \textit{et al.} proposed a
three-layer Huygens metalens operating from 60 to 74~GHz
\cite{Bagheri2023}. The structure used a $50\times50$ array of
two-state unit cells and was illuminated by a dipole antenna. The
maximum reported antenna gain was approximately 25~dBi at 69~GHz.

The Huygens configuration combines electric and magnetic responses to
improve electromagnetic transmission while controlling the
transmitted phase. However, using only two phase states produces
relatively coarse spatial phase quantization compared with a
continuous or multistate phase distribution.

Torres \textit{et al.} experimentally demonstrated a metallic
epsilon-near-zero (ENZ) lens around 144~GHz
\cite{Torres2015}. The structure consisted of narrow metallic
waveguides operating close to their cutoff frequency. The lens
provided a measured transmission enhancement of 15.9~dB and a
directivity of 17.6~dBi.

Although the study confirmed epsilon-near-zero focusing at
millimeter-wave frequencies, its curved metallic construction differs
substantially from low-profile printed metasurface lenses and may be
less suitable for direct integration with planar focal-plane antenna
arrays.


\begin{table}[t]
\centering
\caption{Representative antenna-lens studies}
\label{tab:existing_lens_studies}
\footnotesize
\renewcommand{\arraystretch}{1.15}
\begin{tabularx}{\textwidth}
{
>{\raggedright\arraybackslash}p{2.4cm}
>{\centering\arraybackslash}p{1.6cm}
>{\raggedright\arraybackslash}p{2.1cm}
>{\raggedright\arraybackslash}X
>{\raggedright\arraybackslash}X
}
\toprule
Reference &
Frequency &
Configuration &
Main contribution &
Relevant limitation \\
\midrule

Al-Joumayly and Behdad \cite{Al-Joumayly2011} &
X-band &
MEFSS planar lens &
Wideband subwavelength spatial phase shifters with angular stability &
Multilayer construction; antenna-level evaluation only \\

Yoon and Oh \cite{Yoon2019} &
28~GHz &
Polarization-dependent thin lens &
Polarization-selective focusing with compact unit cells &
Limited transmission-phase range for the focusing polarization \\

Tiwari \textit{et al.} \cite{Tiwari2023} &
28~GHz &
$25\times25$ digital lens &
Discrete beam steering using switchable feed antennas &
Finite beam directions determined by feed positions \\

Lee \textit{et al.} \cite{Lee2022Large} &
28~GHz &
$28\times28$ metamaterial lens &
Large-aperture multibeam gain and system-throughput evaluation &
Limited channel-decorrelation and spatial-mode analysis \\

Lee \textit{et al.} \cite{Lee2021} &
38~GHz &
$13\times13$ single-layer lens &
Broadband and incident-angle-stable focusing &
Scalability toward substantially larger apertures was not demonstrated \\

Bagheri \textit{et al.} \cite{Bagheri2023} &
69~GHz &
$50\times50$ Huygens metalens &
High-gain broadband metalens using electric and magnetic responses &
Coarse two-state phase quantization \\

Torres \textit{et al.} \cite{Torres2015} &
144~GHz &
Metallic ENZ lens &
Experimental ENZ focusing and high directivity &
Nonplanar metallic construction \\

Zeng and Zhang \cite{Zeng2016} &
mmWave &
Analytical lens array &
Path-division multiplexing and angle-dependent focusing model &
Idealized analytical lens response \\

Zeng \textit{et al.} \cite{Zeng2018} &
mmWave &
Full-dimensional lens array &
Multiuser path-based access and delay compensation &
Limited inclusion of physical metasurface imperfections \\

Shen \textit{et al.} \cite{Shen2019} &
28~GHz &
Wideband beamspace MIMO &
Beam-squint-aware beam selection and precoding &
Primarily based on an analytical lens response \\

\bottomrule
\end{tabularx}
\end{table}


\subsection{Research Gaps and Limitations}
\label{subsec:research_gaps}

The reviewed literature demonstrates considerable progress in both
physical antenna-lens design and lens-assisted communication theory.
Nevertheless, several limitations remain relevant to the development
of a practical 38-GHz lens-assisted MIMO system.

First, physical metasurface-lens studies and communication-oriented
lens-MIMO studies have largely developed as separate research
directions. Physical studies commonly characterize unit-cell
transmission, focal-field intensity, gain enhancement, bandwidth,
polarization response, and radiation patterns. In contrast,
communication studies commonly evaluate beamspace sparsity,
RF-chain reduction, beam selection, spectral efficiency, and
path-based multiplexing.

Consequently, the relationship between the electromagnetic focusing
behavior of a practical metasurface lens and the resulting MIMO
channel matrix remains insufficiently investigated. Antenna-gain
enhancement alone does not establish an improvement in spatial
multiplexing. A lens may increase the received signal power while the
channel remains dominated by one spatial mode.

Therefore, a complete evaluation of lens-assisted MIMO should include
not only antenna gain and focal-field intensity but also spatial
correlation, singular-value distribution, condition number, effective
rank, and achievable channel capacity.

Second, although a metasurface lens has already been demonstrated
experimentally around 38~GHz, the scalability of its unit-cell library
toward electrically larger apertures was not established
\cite{Lee2021}. The reported structure used a
$13\times13$ array with a physical aperture of
$39\times39$~mm$^2$.

Therefore, the research gap is not the complete absence of an antenna
lens at 38~GHz. Instead, the limitation is that the demonstrated
38-GHz unit-cell design was evaluated only for a relatively compact
aperture, and its suitability for substantially larger apertures was
not systematically investigated.

A sufficiently wide transmission-phase range allows a unit-cell
library to reproduce the phase distribution required over a large
surface. However, phase coverage alone does not guarantee that the
same unit-cell designs will maintain their performance when arranged
over a substantially larger aperture.

When the aperture is enlarged while the focal distance remains
relatively short, unit cells near the outer regions of the lens are
illuminated at larger local incident angles. Their transmission
magnitude and phase may therefore differ from those obtained using
normal-incidence periodic simulations.

A larger aperture can also require more rapid spatial transitions
between neighboring unit-cell geometries near the aperture edges.
These abrupt transitions may increase spatial discretization errors,
mutual coupling between dissimilar unit cells, and deviations from the
locally periodic approximation.

Most unit-cell libraries are generated by simulating one geometry
under periodic boundary conditions. This simulation assumes that the
unit cell is surrounded by identical copies of itself. In an actual
phase-gradient metasurface lens, however, neighboring unit cells
often have different geometrical dimensions and different
transmission phases.

As a result, the electromagnetic response of a unit cell embedded in
the complete lens may not be identical to the response predicted by
its isolated periodic model. This difference may become more
significant for larger apertures containing sharper spatial phase
variations.

A scalable metasurface-lens unit-cell library should consequently
provide a wide phase-control range, high and relatively uniform
transmission, stable behavior under oblique incidence, sufficiently
fine spatial sampling, low sensitivity to neighboring dissimilar
cells, and acceptable fabrication tolerance.

These characteristics have not been systematically demonstrated for
an electrically larger realization of the selected multilayer
transmitarray topology at 38~GHz. Large-aperture metasurface lenses
have been investigated at other frequencies, including 28~GHz and
approximately 69~GHz \cite{Lee2022Large,Bagheri2023}. However, their
unit-cell structures, material configurations, layer counts, and
operating frequencies differ from the design evaluated in this
thesis. Their results therefore do not directly establish that the
selected 38-GHz unit-cell library can be scaled to a larger aperture
without changing its transmission and focusing behavior.

Third, limited attention has been given to the placement of multiple
receiving antennas within the physical focal region of a practical
metasurface lens. Analytical lens-MIMO models generally assume that
signals arriving from different directions are mapped to clearly
defined antenna locations.

In a practical lens, however, the focal response has a finite spatial
width. The energy associated with one incident direction may therefore
be received by several neighboring antenna elements. Focal-spot
overlap, sidelobes, insertion loss, phase errors, and finite-aperture
diffraction can reduce the separation between different focal-plane
responses.

Consequently, the number, spacing, and positions of the focal-plane
antennas should be selected based on the actual electromagnetic
response of the lens rather than only on an idealized focal-arc
model. These placement parameters may directly affect channel
correlation and the number of usable spatial modes.

Fourth, systematic comparisons between lens-assisted focal-plane
arrays and conventional coplanar arrays remain limited. A fair
comparison requires consistent transmitter positions, propagation
environments, transmit power, bandwidth, antenna count, channel
normalization, and noise assumptions.

Without consistent comparison conditions, an observed improvement may
be caused primarily by increased antenna gain rather than a genuine
improvement in the spatial structure of the MIMO channel.

Accordingly, this thesis investigates a MIMO system assisted by a
metasurface transmitarray lens operating at 38~GHz. It evaluates a
$20\times20$-element aperture designed for focal-plane MIMO reception.

The selected unit-cell topology is implemented in the $20\times20$
array to evaluate the feasibility and electromagnetic performance of
this larger simulated aperture. This single-aperture evaluation does
not, by itself, establish general aperture scalability. The resulting
electromagnetic response is subsequently incorporated into the MIMO
channel analysis.

The focal-plane receiver configuration is compared with a
conventional coplanar antenna arrangement. Instead of evaluating only
gain enhancement or focal-field intensity, the resulting MIMO channel
is characterized using spatial correlation, singular values,
condition number, effective rank, raw link capacity, and normalized
MIMO capacity.

The objective is to determine whether the enlarged physical
transmitarray can provide distinct angle-dependent receive responses
and under which receiver placement, trajectory, aggregation, and
array-order conditions those responses translate into link or
spatial-multiplexing benefits.

% \section{Literature Review}
% \begin{itemize}
%     \item Conventional millimeter-wave MIMO systems
%     \begin{itemize}
%         \item Use of MIMO
%     \end{itemize}
%     \item Antenna lens technologies for beam focusing
%     \item Beamspace MIMO and lens-assisted communication systems
%     \begin{itemize}
%         \item Reference relevant prior work (e.g., studies from Lajos Hanzo's group)
%     \end{itemize}
%     \item Existing research on antenna lenses
%     \begin{itemize}
%         \item Dielectric lens approaches
%         \item Power improvement via aperture advantage
%     \end{itemize}
%     \item Research gaps and limitations
%     \begin{itemize}
%         \item Existing work is largely theoretical, with antenna behavior captured only through mathematical models
%         \item Limited Fidelity Simulation on MIMO System with Antenna lens on FR2 Band e.g., 38 GHz
%     \end{itemize}
% \end{itemize}

% % \paragraph{Propagational Challenges and Hardware Complexity in 5G FR2}
% The utilization of the high-frequency FR2 spectrum is inherently linked to significant propagation challenges arising from its short wavelengths. Signals at these frequencies suffer from severe free-space path loss, substantial penetration losses, and extreme sensitivity to physical blockages, which collectively restrict overall coverage \cite{Zeng2016}. Consequently, maintaining reliable wireless links in FR2 environments requires sophisticated techniques such as directional transmissions and high antenna gains. While the small wavelength size facilitates the integration of numerous antenna elements into a compact aperture, it also necessitates advanced spatial processing capabilities. Conventional approaches to achieve this include Massive MIMO systems designed to provide sufficient beamforming gain; however, fully digital MIMO architectures—which require an independent RF chain for every antenna element—pose substantial practical hurdles at FR2 frequencies. The high-frequency environment demands complex and costly components such as power amplifiers (PAs), low-noise amplifiers (LNAs), mixers, local oscillators, and high-speed ADCs or DACs \cite{Heath2016}. Scaling the number of RF chains in these systems significantly escalates hardware complexity, power consumption, manufacturing costs, and the stringency of calibration requirements.

% % \paragraph{Lens Antenna Array for Beamspace Communication in FR2 MIMO}
% The lens antenna array presents a compelling alternative to conventional MIMO systems by utilizing an electromagnetic lens in conjunction with antenna elements positioned on the focal region or surface. This configuration allows the lens to provide varying phase delays across the aperture, facilitating the passive spatial transformation of signals from the antenna space into a beamspace \cite{shen2019}. By focusing waves originating from different directions onto specific subsets of antenna elements, only those receiving dominant energy need to be connected to the RF chains. This characteristic is particularly advantageous for Frequency Range 2 (FR2) MIMO systems, where propagation channels are typically characterized by a limited number of dominant paths and a relatively sparse beamspace \cite{Zeng2016}. In such environments, most channel energy is concentrated on only a few beams or antennas, allowing for efficient beam selection to significantly reduce the required number of RF chains. Furthermore, the integration of lens antennas can enhance system performance by improving angular separation, spatial channel decorrelation, effective channel rank, spatial multiplexing capability, and overall channel capacity \cite{Zeng2016}.

% % \paragraph{Metasurface-Based Planar Lens Solutions for High-Frequency Communication}
% While conventional dielectric lenses are valued for their high focusing capabilities and gain, they possess several significant drawbacks that hinder their integration into modern wireless systems. These include thick profiles, substantial weight, curved surfaces, and reflection losses stemming from impedance mismatches \cite{Al-Joumayly2011}. Furthermore, dielectric lenses can be bulky, expensive to manufacture at certain microwave frequencies, and inherently difficult to integrate with planar antenna arrays due to their complex geometries. To overcome these limitations, there is a growing shift toward planar lens architectures. A prominent solution in this area is the metasurface-based planar antenna lens, which comprises an array of sub-wavelength unit cells that act as individual spatial phase shifters \cite{Al-Joumayly2011}. By adjusting the geometric dimensions of these unit cells, one can precisely control the transmission magnitude, phase, and frequency response across the aperture. The distribution of these cells allows for the synthesis of specific phase profiles capable of transforming spherical wavefronts into planar ones or vice versa, effectively focusing waves while maintaining a compact and low-profile form factor \cite{Al-Joumayly2011}.

% \section{Research Motivation}
% \begin{itemize}
%     \item Channel correlation in conventional coplanar MIMO arrays
%     \item Potential advantages of focal-plane antenna placement
%     \item Need for a systematic evaluation of antenna placement effects on MIMO performance
% \end{itemize}

\section{Motivation}
% \paragraph{Design Motivation and Requirements for 5G FR2 Lens Antennas}
The primary motivation for this research is the need for compact front ends that can be integrated into 5G FR2 systems. Conventional dielectric lenses can be thick, heavy, curved, and difficult to integrate with planar arrays. By contrast, a planar transmitarray assembled from printed unit cells can provide spatial phase control in a low-profile form and can be combined with patch antennas positioned near its focal region. At 38~GHz, passive focusing may strengthen a direction-matched link without requiring an active phase shifter at every aperture element, while an incidence-angle-dependent focal response may support multiple focal-plane receiver (RX) branches.

% Poin utama Paragraf 1:
% Mengembangkan lens yang planar dan low-profile.
% Menghindari struktur dielectric lens yang tebal dan melengkung.
% Menggunakan teknik fabrikasi PCB.
% Mengintegrasikan lens dengan patch antenna dan focal-plane receiver.
% Menghasilkan phase coverage yang besar dengan insertion loss yang tetap terkendali.

% Poin utama Paragraf 2:
% Metasurface lens mengontrol phase gelombang tanpa active phase shifter.
% Fokus yang dihasilkan dapat meningkatkan received power.
% Fokus berubah sesuai incidence angle.
% Beberapa antenna receiver dapat ditempatkan pada focal arc atau focal plane.

% \paragraph{Beyond Antenna Gain Enhancement}
This thesis deliberately evaluates effects beyond antenna gain. The central question is whether the transmitarray alters the channel matrix by separating direction-dependent fields in the focal region. Such a transformation may lower spatial correlation, balance the singular values, reduce the condition number, and increase effective rank and MIMO capacity. These outcomes are treated as hypotheses to be tested at multiple scales rather than as assumed consequences of focusing. Chapters~2--4 show why this distinction is necessary: HFSS SBR+ and Hybrid indicate improved normalized $3\times3$ mode balance, whereas Sionna-RT shows a much smaller component-level change. Moreover, the paired mobile $3\times3$ receiver-configuration comparison does not show an aggregate capacity or effective-rank advantage under equal weighting across the three RX branches. Instead, the lens provides strong local angular selectivity whose utility depends on alignment and receiver operation.

\section{Research Questions}

The investigation is organized around four research questions:
\begin{enumerate}
    \item Can the selected 38~GHz unit-cell library be implemented as a simulated $20\times20$-element transmitarray that produces a usable focal region and angle-dependent receive patterns?
    \item Which lens-induced channel trends are consistent across HFSS SBR+, HFSS Hybrid, and Sionna-RT when raw link gain is separated from normalized spatial-mode structure?
    \item How do the complete $3\times3$ focal-plane lens-assisted and coplanar without-lens RX designs compare along matched linear and half-circular transmitter (TX) trajectories in terms of link-budget capacity, conditioning, effective rank, normalized constructed-MIMO capacity, and spatial decorrelation?
    \item How do total normalized constructed-MIMO capacity and per-dimension metrics differ among the complete lens-assisted $3\times3$, $5\times5$, and $9\times9$ configurations along a common trajectory?
\end{enumerate}

% Poin utama:
% Lens bukan hanya gain enhancer.
% Lens dapat mengubah struktur channel matrix.
% Pemisahan field pada focal region dapat menurunkan channel overlap.
% Hipotesis penelitian:
% spatial correlation menurun;
% singular values menjadi lebih seimbang;
% condition number membaik;
% effective rank meningkat;
% channel capacity meningkat.

% \section{Research Objectives}
% \begin{itemize}
%     \item Develop a 38 GHz antenna lens-assisted MIMO system
%     \item Investigate the influence of antenna placement on MIMO channel decorrelation
%     \item Evaluate spatial multiplexing capability using different antenna configurations
%     \item Validate the proposed channel-modeling approach
%     \item Demonstrate the communication benefits of antenna lens-assisted MIMO systems
% \end{itemize}

\section{Research Objectives}
Building on this motivation, the specific objectives of this thesis are as follows:
\begin{itemize}
    \item To design and characterize a 38~GHz patch feed antenna and a planar $20\times20$-element metasurface transmitarray, including its simulated focal distance, focal-arc response, and angle-dependent radiation patterns in Ansys HFSS\@.
    \item To formulate a common framework for analyzing HFSS S-parameters and Sionna-RT channel responses using raw link metrics, unit-Frobenius normalization at each frequency, and unit-mean-coefficient-power normalization at each waypoint.
    \item To assess whether lens-induced changes are consistent across HFSS SBR+, HFSS Hybrid, and Sionna-RT while separating raw link changes from normalized changes in spatial-mode structure.
    \item To compare complete lens-assisted and without-lens $3\times3$ RX designs at matched waypoints along prescribed linear and half-circular trajectories, including both global and local spatial-decorrelation estimators.
    \item To characterize configuration-level higher-order scaling for complete lens-assisted $3\times3$, $5\times5$, and $9\times9$ systems on the same linear trajectory using total normalized constructed-MIMO capacity and per-dimension metrics.
\end{itemize}

\section{Thesis Organization}

The remainder of this thesis is organized as follows. Chapter~2, \textit{Methodology and Numerical Model Verification}, describes the 38~GHz patch antenna and $20\times20$-element metasurface transmitarray, including the unit-cell response, simulated focal distance, focal-arc behavior, and imported angle-dependent RX patterns. It then defines the channel state information (CSI), singular-value, correlation, condition number, effective rank, and capacity metrics used throughout the thesis. The $3\times3$ lens-assisted and without-lens channels are compared using HFSS SBR+, HFSS Hybrid, and Sionna-RT, with raw and per-frequency unit-Frobenius-normalized results reported separately.

Chapter~3, \textit{MIMO Communication Performance Evaluation}, contains two trajectory-based evaluation studies in Sionna-RT\@. The first is a paired comparison of complete lens-assisted and without-lens $3\times3$ RX designs at identical waypoints along linear and half-circular TX paths. It reports link-budget capacity, constructed-MIMO conditioning, effective rank, fixed-10-dB constructed-MIMO capacity after unit-mean-coefficient-power normalization, and spatial decorrelation at global-pooled, waypoint-local, and block-median aggregation scopes. The second is a configuration-level comparison of complete lens-assisted $3\times3$, $5\times5$, and $9\times9$ systems along the same linear trajectory. Because nominal total TX power, physical aperture, and the RX angle grid vary with order, this study describes joint configuration scaling rather than an order-only effect.

Chapter~4, \textit{Conclusions and Future Work}, synthesizes the electromagnetic, cross-solver, mobile-placement, and higher-order findings. It presents the conditional system-level conclusion, identifies methodological and reproducibility limitations, and proposes experimental validation, practical beam management, factor-controlled placement studies, fair power-and-aperture scaling, statistical convergence analysis, lens re-optimization, richer propagation scenarios, and multiuser extensions.

% ## Chapter 1: Introduction

% ### 1.1 Background

% * Evolution of millimeter-wave communication systems

%   * FR1(Cellular Band maybe become full, move to mmwave)
% * Challenges of line-of-sight-dominant MIMO channels

%   * Share abouts many mmwave challenges, Show Paper Proof
% * Importance of channel decorrelation and spatial multiplexing

%   * idk
% * Motivation for antenna lens-assisted MIMO systems

%   * Should to improve channel decorellation and Spatial MultiPlexing

% ### 1.2 Literature Review

% * Conventional millimeter-wave MIMO systems

%   * Using MIMO
% * Antenna lens technologies for beam focusing
% * Beamspace MIMO and lens-assisted communication systems

%   * Paper muridnya Lajos Hanzo
%   *
% * Existing research on antenna lenses

%   * Using Dielectric Lens
%   * Just in case improve the Power with Advantage Aparture
% * Research gaps and limitations
%   * Commonly Only Based On Theory, Antenna still on Math Model.

% ### 1.3 Research Motivation

% * Channel correlation in conventional coplanar MIMO arrays
% * Potential advantages of focal-plane antenna placement

%   * We try to
% * Need for a systematic evaluation of antenna placement effects on MIMO performance

% ### 1.4 Research Objectives

% * Develop a 38 GHz antenna lens-assisted MIMO system
% * Investigate the influence of antenna placement on MIMO channel decorrelation
% * Evaluate spatial multiplexing capability using different antenna configurations
% * Validate the proposed channel-modeling approach
% * Demonstrate the communication benefits of antenna lens-assisted MIMO systems

% ### 1.5 Thesis Organization
